% Introduction
\section{Introduction}

\subsection{Background: Code Generation and the Learning Signal Problem}

Code generation presents unique challenges for language model training. Unlike natural language, code exhibits several distinctive characteristics:

\begin{itemize}[nosep]
    \item \textbf{Long sequences}: Competitive programming solutions often exceed 1,000 tokens
    \item \textbf{Sparse, episodic rewards}: A solution is either correct or incorrect---no partial credit
    \item \textbf{Critical token sensitivity}: Function names, conditional expressions, and return statements have disproportionate impact on correctness
\end{itemize}

Standard supervised fine-tuning (SFT) with cross-entropy loss treats all tokens equally, allocating identical learning signal to boilerplate code and critical logic. Multi-Token Prediction (MTP) \citep{gloeckle2024better} extends this by predicting multiple future tokens simultaneously, improving training efficiency and representation quality. However, MTP inherits the same limitation: \textbf{uniform loss weighting fails to reflect token importance differences}.

This observation motivates our core hypothesis: \textit{``Not All Tokens Are What You Need''}---allocating more learning resources to important tokens should improve performance and convergence efficiency at equivalent FLOPs.

\subsection{Research Questions}

We formulate three research questions to investigate token-weighted learning for code generation:

\begin{description}[style=unboxed,leftmargin=0cm]
    \item[RQ1:] Under what conditions does GAE-weighted cross-entropy improve code generation performance?
    \begin{itemize}[nosep]
        \item \textbf{RQ1a:} What is the effect on out-of-domain benchmarks (short function-level tasks like HumanEval/MBPP)?
        \item \textbf{RQ1b:} What is the effect on in-domain benchmarks (long competitive programming like CodeContests)?
    \end{itemize}

    \item[RQ2:] How can we stably learn token-level value estimates?
    \begin{itemize}[nosep]
        \item Pairwise ranking vs. MSE regression
        \item Length-balanced stratified sampling
    \end{itemize}

    \item[RQ3:] How does critic quality affect downstream policy performance? (Exploratory, incomplete)
\end{description}

\subsection{Proposed Approach: Two-Phase WMTP}

We propose \textbf{Weighted Multi-Token Prediction (WMTP)}, a two-phase pipeline illustrated in Figure~\ref{fig:pipeline}:

\begin{description}[style=unboxed,leftmargin=0cm]
    \item[Phase 1 -- Critic Training:] Train an independent value model on correct/incorrect solution pairs using pairwise ranking loss. The critic learns token-level value estimates $V_t$ that distinguish solutions likely to be correct from those likely to be incorrect.

    \item[Phase 2 -- Weighted Policy Training:] Freeze the critic and use it to compute token advantages via GAE. Apply whitening normalization with EMA running statistics, convert to exponential weights, and modulate the MTP cross-entropy loss.
\end{description}

% Placeholder for pipeline figure
\begin{figure}[t]
    \centering
    \fbox{\parbox{0.9\textwidth}{
        \centering
        \textbf{[Figure 1: WMTP Pipeline]}\\[1em]
        Phase 1: CodeContests (correct/incorrect pairs) $\rightarrow$ Pairwise Ranking $\rightarrow$ Critic $V_t$\\[0.5em]
        Phase 2: Critic frozen $\rightarrow$ GAE $A_t$ $\rightarrow$ Whitening + EMA $\rightarrow$ $w_t = \exp(\hat{A}_t/\beta)$ $\rightarrow$ Weighted MTP Loss
    }}
    \caption{Overview of the WMTP two-phase pipeline. Phase 1 trains an independent critic using pairwise ranking on solution pairs. Phase 2 uses the frozen critic to compute token-level advantages for weighting the MTP loss. The critic is \textbf{not used at inference time}.}
    \label{fig:pipeline}
\end{figure}

\subsection{Contributions}

Our main contributions are:

\begin{enumerate}[nosep]
    \item \textbf{GAE-weighted CE exploration for LLM code generation}: We transfer GAE-based advantage estimation, primarily used in robotics control, to token weighting for code generation, observing positive signals on out-of-domain benchmarks.

    \item \textbf{Pairwise ranking for value learning}: We replace pointwise MSE regression with Bradley-Terry pairwise ranking, providing stable learning signals appropriate for the code domain where relative correctness is easier to judge than absolute quality.

    \item \textbf{Sparse reward stabilization recipe}: We combine whitening, EMA normalization, exponential weighting, and clipping to control learning instability from sparse episodic rewards.

    \item \textbf{Out-of-domain transfer observation}: Using only 1.4\% of available training data (50K/3.6M samples), we observe improvements on HumanEval, MBPP, and GSM8K (causal relationship not confirmed).

    \item \textbf{Honest limitation disclosure}: We explicitly acknowledge unverified MTP necessity (vs. NTP), 63\% critic accuracy ceiling, in-domain performance decline, and absence of random weight baselines.
\end{enumerate}

\textbf{Important note:} MTP is an \textit{implementation choice} in this work. The necessity of MTP over standard NTP for token weighting effectiveness requires separate verification, which we leave to future work.
